\documentclass[a4paper,12pt]{article} % Type de document
\usepackage[utf8]{inputenc} % Pour génerer les accents automatiquement
\usepackage[T1]{fontenc} % Pour gerér la police des accents
\usepackage[french]{babel} % Pour la typographie française
%Ce sont les trois packages utiles pour écrire en français
\frenchbsetup{StandardLists=true}%afin d'éviter tout conflit entre l'extension %\usepackage[french]{babel}.

\parindent=0cm %pour le changement d'indentation d'un paragraghe
%\usepackage{lipsum}pour generer du faux texte

\usepackage{amsmath} % Pour écrire des symboles mathématiques
\usepackage{amsfonts} % Pour gerer les doubles barres
\usepackage{amsthm}
\usepackage{enumitem}% Pour personnaliser des listes
\usepackage{amssymb} % Pour utiliser plus de symboles 
\usepackage{xcolor}
%\usepackage[margin=2.5cm]{geometry}% Pour regler les marges indentiquement
\usepackage[left=2cm,top=2cm,bottom=2.7cm,right=2cm]{geometry}% Pour regler les marges differemment
%\usepackage{setspace}\renewcommand{\baselinestretch}{1.5}%Pour modifier l'interligne
\usepackage{systeme}
\usepackage{xstring}

\usepackage{hyperref}
\newtheorem{theorem}{Théorème}[section]
\newtheorem{corollary}{Corollaire}[theorem]
\newtheorem{lemma}{Lemme}[section]
\newtheorem{proposition}{Proposition}[section]
\newtheorem{property}{Propriété}[section]
\newtheorem{definition}{Définition}[section]
\newtheorem*{remark}{Remarque}% Pour ajouter des remarques non numérotées
\newtheorem*{example}{Exemple} % Pour ajouter des exemples non numérotés
\newtheorem*{preuve}{\textbf{Preuve}}

%\newenvironment{preuve}[1][Preuve]{\textbf{#1.} }{\ \rule{0.5em}{0.5em}}

%\renewcommand{\arraystretch}{2.2}

%\title{\textbf{Sujet: Modèle de type Bloch avec les termes de relaxation de Pauli et l'interaction de Coulomb.}} % Titre du document
%\author{\textbf{OKAINGNI OKPO DORGELES} \\
%\date{}\\ %Date de création
%\textbf{Étudiant en Master 2 Maths-Info} \\
%\textbf{CFC de DALOA}}

\begin{document}
%\maketitle{} % Pour generer le titre

\tableofcontents

\bigskip



\section{Modèle de type Bloch avec les termes de relaxation, de Coulomb et le champ électromagnétique et la matière}

La forme générale du modèle de type Bloch avec les termes de relaxation, de Coulomb et le champ électromagnétique et la matière est donnée par :
\begin{equation}\label{modèle de base}
\left\{
   \begin{array}{rcl}
   
   \partial_{t}\rho & = & -\dfrac{i}{\hbar}\Big([E_{0},\rho]+[V^{\textrm{C}}(\rho),\rho]+E(R,t)\cdot[\textbf{\textrm{m}},\rho]\Big)+ W(\rho), \forall t>0 \\
   \rho(0)          & = & \rho^{0},
   
   \end{array}
 \right.
\end{equation}

où 
\begin{itemize}[label=$\centerdot$]
\item $ \rho^{0}\in \mathcal{M}_{N}(\mathbb{C}) $,
\item $ E_{0}= diag((e_{j})_{j}) $  où $j\in \{1,\ldots, N\}, \; e_{j}$ est une énergie propre libre du niveau quantique $j$ et représente une valeur propre de l'opérateur de Schrödinger non stationnaire et sans les potentiels d'interaction,
\item $ W(\rho)\in \mathcal{M}_{N}(\mathbb{C}) $ correspond aux termes de relaxation issus de l'équation maîtresse de Pauli donnés dans \cite{mareference2},
\item $ V^{\textrm{C}}(\rho)\in \mathcal{M}_{N}(\mathbb{C}) $ correspond aux termes de Coulomb issus de l'interaction de Coulomb entre les particules (électrons et trous) donnés dans \cite{mareference1}. Les termes de Coulomb correspondent au comportement des particules dans les systèmes quantiques à plusieurs particules. Ces particules étant chargées se repoussent ou s’attirent et ce sont ces deux actions (répulsion et attraction) qui constituent l’interaction de Coulomb (voir dans \cite{mareference1} à la page 31). 
\item $\textbf{\textrm{M}}(R,t)=E(R,t)\cdot \textbf{\textrm{m}} $ avec $ \textbf{\textrm{m}}$ la matrice représentant le moment dipolaire, $ E(R,t)$ le champ électromagnétique et le vecteur $ R $  est la position macroscopique de la boite considérée (voir dans \cite{mareference1}, pages 27-28).                         
 \end{itemize}

\section{Le système dynamique réel associé au modèle de type Bloch avec les termes de relaxation et l'interaction de Coulomb}

Le modèle de type Bloch avec les termes de relaxation, de Coulomb est donnée par :
\begin{equation}\label{modèle de base}
\left\{
   \begin{array}{rcl}
   
   \partial_{t}\rho & = & -\dfrac{i}{\hbar}\Big([E_{0},\rho]+[V^{\textrm{C}}(\rho),\rho]\Big)+ W(\rho), \forall t>0 \\
   \rho(0)          & = & \rho^{0},
   
   \end{array}
 \right.
\end{equation}
Le système dynamique réel associé à ce modèle de type Bloch est .
\begin{equation}\label{equation du systeme dynamique reel non lineaire}
\partial_{t}U(t) = A(t,U(t))U(t) 
\end{equation}
On a: $ A(.,U)= A^{L}+ A^{N}(.,U)$ \\
où 
\begin{equation}
A^{L}=
\begin{pmatrix}
W_{1} &      0                    & 0                       \\
    0 & D_{\gamma}                &  -\big(\mathcal{E}_{0}+\mathcal{R}^{0}_{C}\big)  \\
    0 & \mathcal{E}_{0}+ \mathcal{R}^{0}_{C}  &  D_{\gamma}
\end{pmatrix}\label{matrice Al}
\end{equation}

$$
 W_{1}= 
\begin{pmatrix}
-(w_{21} + w_{31}) & w_{21}\dfrac{a_{2}}{a_{1}} & w_{31}\dfrac{a_{3}}{a_{1}}\\
            w_{21} & -(w_{21}\dfrac{a_{2}}{a_{1}} + w_{32}) & w_{32}\dfrac{a_{3}}{a_{2}}\\
            w_{31} &                                 w_{32} & -(w_{31}\dfrac{a_{3}}{a_{1}} + w_{32}\dfrac{a_{3}}{a_{2}})
\end{pmatrix}
$$
et 
$$ D_{\gamma} = \textrm{diag}(\gamma_{12},\gamma_{13}, \gamma_{23}), \quad \mathcal{E}_{0} = - \dfrac{1}{\hbar}  \textrm{diag}(e_{1}-e_{2},e_{1}-e_{3}, e_{2}-e_{3}) $$

et 
$$ \mathcal{R}^{0}_{C} =- \dfrac{1}{\hbar}  \textrm{diag}(0,\nu^{C}, \nu^{C})$$
où \quad
  $ \nu^{C} = 2(R^{\textsc{v}}_{2112}-R^{\textsc{v}}_{2121})+R^{\textsc{c-v}}_{3131} + R^{\textsc{c-v}}_{3232} $ , $ \nu^{C} \in \mathbb{R}. $\\
\begin{equation}\label{matrice AN}
A^{N}(.,U)=
\begin{pmatrix}
0 & \mathcal{R}^{1}_{C}(U) & \mathcal{R}^{2}_{C}(U)\\
\mathcal{R}^{3}_{C}(U) & \mathcal{R}^{4}_{C}(U) & \mathcal{R}^{5}_{C}(U) \\
\mathcal{R}^{6}_{C}(U) & \mathcal{R}^{7}_{C}(U)&\mathcal{R}^{8}_{C}(U)
\end{pmatrix}
\end{equation}
 où
  \begin{eqnarray*}
 \begin{split}
\mathcal{R}^{1}_{C}(U) &= \dfrac{2}{\hbar}
\begin{pmatrix}
\Im(R^{\textsc{c-v}}_{3132})U_{3}& -\Im(R^{\textsc{c-v}}_{3132})U_{6} & \Re(R^{\textsc{c-v}}_{3132})U_{8}\\
-\Im(R^{\textsc{c-v}}_{3132})U_{3}& \Im(R^{\textsc{c-v}}_{3132})U_{6}& -\Re(R^{\textsc{c-v}}_{3132})U_{8}\\
            0 &                               0 & 0
\end{pmatrix} \quad \textrm{et} \\
\mathcal{R}^{2}_{C}(U) &= \dfrac{2}{\hbar}
\begin{pmatrix}
-\Re(R^{\textsc{c-v}}_{3132})U_{3} & -\Im(R^{\textsc{c-v}}_{3132})U_{9} & -\Re(R^{\textsc{c-v}}_{3132})U_{5}\\
\Re(R^{\textsc{c-v}}_{3132})U_{3} & \Im(R^{\textsc{c-v}}_{3132})U_{9} & \Re(R^{\textsc{c-v}}_{3132})U_{5}\\
            0 &                               0 & 0
\end{pmatrix}
\end{split}
\end{eqnarray*}

 \begin{eqnarray*}
\begin{split}
\mathcal{R}^{3}_{C}(U)&= \dfrac{1}{\hbar}
\begin{pmatrix}
-\Im(R^{\textsc{c-v}}_{3132})U_{3} & \Im(R^{\textsc{c-v}}_{3132})U_{3} & 0\\
\Im(R^{\textsc{c-v}}_{3132})U_{6}+\Re(R^{\textsc{c-v}}_{3132})U_{9} &
-\big[2(R^{\textsc{v}}_{1221}-R^{\textsc{v}}_{1212})+R^{\textsc{c-v}}_{3232}\big]U_{8} & 0 \\
-\big[2(R^{\textsc{v}}_{1221}-R^{\textsc{v}}_{1212})+R^{\textsc{c-v}}_{3131}\big]U_{9} & 
\Re(R^{\textsc{c-v}}_{3132})U_{8}- \Im(R^{\textsc{c-v}}_{3132})U_{5} & 0
\end{pmatrix}\\
\end{split}
\end{eqnarray*}

\begin{small}
\begin{eqnarray*}
\begin{split}
\mathcal{R}^{4}_{C}(U)&= \dfrac{1}{\hbar}
\begin{pmatrix}
0 & \Im(R^{\textsc{c-v}}_{3132})U_{5} & -\Im(R^{\textsc{c-v}}_{3132})U_{6} \\
-\Re(R^{\textsc{c-v}}_{3132})U_{8} & -\Im(R^{\textsc{c-v}}_{3132})U_{4} & \big[2(R^{\textsc{v}}_{1221}-R^{\textsc{v}}_{1212})+R^{\textsc{c-v}}_{3232}\big]U_{7}\\
\Im(R^{\textsc{c-v}}_{3132})U_{6} & -\big[2(R^{\textsc{v}}_{1221}-R^{\textsc{v}}_{1212})+R^{\textsc{c-v}}_{3131}\big]U_{7} & -\Re(R^{\textsc{c-v}}_{3132})U_{7}
\end{pmatrix}
\end{split}
\end{eqnarray*}


\begin{eqnarray*}
\begin{split}
\mathcal{R}^{5}_{C}(U)&= \dfrac{1}{\hbar}
\begin{pmatrix}
U_{3} & \Im(R^{\textsc{c-v}}_{3132})U_{8}-U_{6} & U_{5}-\Im(R^{\textsc{c-v}}_{3132})U_{9} \\
\Re(R^{\textsc{c-v}}_{3132})U_{5} & -\Im(R^{\textsc{c-v}}_{3132})U_{7} & \big[2(R^{\textsc{v}}_{1221}-R^{\textsc{v}}_{1212})+R^{\textsc{c-v}}_{3232}\big]U_{4}\\
-\Im(R^{\textsc{c-v}}_{3132})U_{9} & \big[2(R^{\textsc{v}}_{1221}-R^{\textsc{v}}_{1212})+R^{\textsc{c-v}}_{3131}\big]U_{4} & -\Re(R^{\textsc{c-v}}_{3132})U_{4}
\end{pmatrix}
\end{split}
\end{eqnarray*}
\end{small}

\begin{eqnarray*}
\begin{split}
\mathcal{R}^{6}_{C}(U)&= \dfrac{1}{\hbar}
\begin{pmatrix}
\Re(R^{\textsc{c-v}}_{3132})U_{3} & -\Re(R^{\textsc{c-v}}_{3132})U_{3} & 0 \\
-\Re(R^{\textsc{c-v}}_{3132})U_{6}+\Im(R^{\textsc{c-v}}_{3132})U_{9} & \big[2(R^{\textsc{v}}_{1221}-R^{\textsc{v}}_{1212})+R^{\textsc{c-v}}_{3232}\big]U_{5} & 0 \\
\big[2(R^{\textsc{v}}_{1221}-R^{\textsc{v}}_{1212})+R^{\textsc{c-v}}_{3131}\big]U_{6} & -\Re(R^{\textsc{c-v}}_{3132})U_{5}- \Im(R^{\textsc{c-v}}_{3132})U_{8} & 0 
\end{pmatrix}\\
\mathcal{R}^{7}_{C}(U)&= \dfrac{1}{\hbar}
\Bigg(\begin{smallmatrix}
-(R^{\textsc{c-v}}_{3131}-R^{\textsc{c-v}}_{3232})U_{3} & -\Re(R^{\textsc{c-v}}_{3132})U_{5}
 +(R^{\textsc{c-v}}_{3131}-R^{\textsc{c-v}}_{3232})U_{6} &  \Re(R^{\textsc{c-v}}_{3132})U_{6}\\
\Im(R^{\textsc{c-v}}_{3132})U_{8} & \Re(R^{\textsc{c-v}}_{3132})U_{4} & -\big[2(R^{\textsc{v}}_{1221}-R^{\textsc{v}}_{1212})+R^{\textsc{c-v}}_{3232}\big]U_{4}\\
\Im(R^{\textsc{c-v}}_{3132})U_{9} & -\big[2(R^{\textsc{v}}_{1221}-R^{\textsc{v}}_{1212})+R^{\textsc{c-v}}_{3131}\big]U_{4} & \Re(R^{\textsc{c-v}}_{3132})U_{4} 
\end{smallmatrix}\Bigg)
\end{split}
\end{eqnarray*}

\begin{small}
\begin{eqnarray*}
\begin{split}
\mathcal{R}^{8}_{C}(U)&= \dfrac{1}{\hbar}
\begin{pmatrix}
0 & -\Re(R^{\textsc{c-v}}_{3132})U_{8} & \Re(R^{\textsc{c-v}}_{3132})U_{9}+(R^{\textsc{c-v}}_{3131}-R^{\textsc{c-v}}_{3232})U_{8}\\
\Im(R^{\textsc{c-v}}_{3132})U_{5} & \Re(R^{\textsc{c-v}}_{3132})U_{7} &  \big[2(R^{\textsc{v}}_{1221}-R^{\textsc{v}}_{1212})+R^{\textsc{c-v}}_{3232}\big]U_{7}\\
\Im(R^{\textsc{c-v}}_{3132})U_{6} & -\big[2(R^{\textsc{v}}_{1221}-R^{\textsc{v}}_{1212})+R^{\textsc{c-v}}_{3131}\big]U_{7} & -\Re(R^{\textsc{c-v}}_{3132})U_{7}
\end{pmatrix}
\end{split}
\end{eqnarray*}
\end{small}



















































\bibliographystyle{plain}
\bibliography{mabibliomaster}
 \addcontentsline{toc}{section}{Bibliographie}
\end{document}